\documentclass{article}
\usepackage[UTF8]{ctex}
\usepackage{amsmath,amssymb,amsfonts}
\usepackage{enumitem}
\usepackage{geometry}
\geometry{a4paper, margin=1in}
% === 错题本排版优化（不参与题目标识） ===
\setlength{\parindent}{0pt}
\setlength{\parskip}{0.35em plus 0.10em minus 0.08em}
\setlength{\emergencystretch}{3em}
\setlist{itemsep=0.25em, topsep=0.40em, parsep=0pt}
\allowdisplaybreaks[2]
% === 错题本排版优化结束 ===
\begin{document}
\section*{华中师范大学2025年高等代数真题}
\section*{130．华中师范大学2025年高等代数试卷}
一、填空题、每题5分，共20分。

\begin{itemize}

\subsection{题目1}
  \item[1．] 设矩阵
\[
A=\left(\begin{array}{ccc}
1 & 0 & 1 \\
0 & 1 & 1 \\
0 & a & -1
\end{array}\right), B=\left(\begin{array}{llll}
2 & 0 & 2 & 4 \\
1 & 2 & 3 & 4 \\
5 & 6 & 7 & 8
\end{array}\right) .
\]
若 \(\operatorname{rank}(A B)<3\) ，求 \(a=\) \(\_\_\_\_\)．

\subsection{题目2}
  \item[2．] 当 \(x=\) \(\_\_\_\_\)时，矩阵 \(\left(\begin{array}{ccc}2 & 0 & 1 \\ 3 & 1 & x \\ 4 & 0 & 5\end{array}\right)\) 可相似对角化、
\end{itemize}

\subsection{题目3}
3．\(n\) 阶实对称矩阵有 \(\_\_\_\_\)种不同的合同标准形。

\subsection{题目4}
4．若3阶复方阵 \(A\) 恰有两个互异的特征值 \(\lambda_{1}=1, \lambda_{2}=-1, A\) 不能相似对角化且 \(\operatorname{tr}(A)=-1\) 。则\(A\) 的初等因子组为 \(\_\_\_\_\)。

二、解答题。\\

\subsection{题目5}
1．（20 分）计算行列式

\[
\left|\begin{array}{cccc}
1^{2023} & 2^{2023} & \cdots & 2025^{2023} \\
2^{2023} & 3^{2023} & \cdots & 2026^{2023} \\
\vdots & \vdots & & \vdots \\
2025^{2023} & 2026^{2023} & \cdots & 4049^{2023}
\end{array}\right|
\]

其中所有元素指数均为2023，各行各列的底数为等差数列。

\subsection{题目6}
2．（25 分）设 \(A, A_{1}, B_{1}\) 为复数域上的 \(n\) 阶方阵。

\begin{itemize}
  \item[（1）] 如果对于任意自然数 \(k \geq 1\) ，有 \(\operatorname{tr}\left(A^{k}\right)=0\) ，证明：\(A\) 是幂零矩阵。
  \item[（2）] 如果 \(A=A_{1} B_{1}-B_{1} A_{1}\) ，且 \(A A_{1}=A_{1} A\) ，证明：\(A\) 是幂零矩阵。
\end{itemize}

\subsection{题目7}
3．（10 分）证明：\(f(x)=x^{3}+x^{2}+2025\) 在有理数域上不可约．

\subsection{题目8}
4．（25 分）设 \(f(X)=X^{\mathrm{T}} A X\) 是正定二次型，\(A=\left(a_{i j}\right)_{n \times n}\) 是实对称阵，其中 \(X=\left(x_{1}, x_{2}, \cdots, x_{n}\right)^{\mathrm{T}}\) 是\(n\) 维实向量．\\
（1）证明：关于 \(Y=\left(y_{1}, y_{2}, \cdots, y_{n}\right)^{\mathrm{T}}\) 的二次型 \(g(Y)=\operatorname{det}\left(\begin{array}{cc}A & Y \\ Y^{\mathrm{T}} & 0\end{array}\right)\) 是负定二次型。\\
（2）证明：\(|A| \leq a_{n n} A_{n-1}\) ，其中 \(A_{n-1}\) 是 \(A\) 的 \(n-1\) 阶顺序主子式。\\
（3）证明：\(|A| \leq a_{11} a_{22} \cdots a_{n n}\) ．\\
（4）对实方阵 \(B=\left(b_{i j}\right)_{n \times n}\) ，证明：\(|B| \leq \sqrt{\prod_{i=1}^{n}\left(b_{1 i}^{2}+b_{2 i}^{2}+\cdots+b_{n i}^{2}\right)}\) ．

\subsection{题目9}
5。（25分）设典型基矩阵 \(E_{i j}\) 是第 \(i\) 行第 \(j\) 列元素为 1 ，其余元素为 0 的 \(n\) 阶方阵。

\begin{itemize}
  \item[（1）] 讨论 \(E_{i j} E_{k l}\) 。
  \item[（2）] 对 \(n\) 阶方阵 \(A, B\) ，证明： \(\operatorname{tr}(A B-B A)=0\) 。
  \item[（3）] 记 \(W\) 是由 \(A B-B A\) 生成的线性子空间，其中 \(A, B\) 取遍 \(n\) 阶方阵。求 \(\operatorname{dim} W\) 。
\end{itemize}

\subsection{题目10}
6．（25 分）设1和2是4阶复方阵 \(A\) 的所有非零特征值， \(\operatorname{rank}(A)=3\) 。

\begin{itemize}
  \item[（1）] 写出 \(A\) 的所有可能的若尔当标准形（不考虑若尔当块的顺序）。
  \item[（2）] 若 \(\operatorname{tr}(A)=5\) 且 \(A\) 不能相似对角化，求 \(A\) 的特征多项式和极小多项式。
  \item[（3）] 在（2）的条件下，对任意的正整数 \(k\) ，求 \(A^{k}\) 的若尔当标准形．
\end{itemize}
\end{document}
